\chapter*{Annexes}
\addcontentsline{toc}{chapter}{Annexes}

\section*{Recours à l'IA générative pour assistance rédactionnelle}

L'intelligence artificielle générative a été utilisée comme appui rédactionnel. Son rôle s'est limité
à la reformulation de phrases, à la recherche de synonymes et à l'ajustement du choix des mots, dans
un objectif de clarté et de fluidité. Le contenu, la structure et les idées relèvent d'un travail personnel
; l'outil n'a servi qu'à optimiser la forme linguistique, sans incidence sur le fond.
Sur l'aspect technique, l'IA générative a servi à corriger de petits bugs dans des scripts, à mettre en
forme des visuels sous LaTeX et à orienter la conception de quelques graphiques à partir de résultats
existants. Ces interventions ont principalement permis de gagner du temps sur la mise en page et
de clarifier certaines étapes du code. Les arbitrages d'analyse, la méthodologie et l'interprétation
demeurent indépendants.

